\begin{lemma}[Successive parallel extensions at one support share a color]
Let \(u\in V(T_{q+1})\).  Suppose that
\[
 b_0<b_1<\cdots<b_p\leq q
\]
is a nonempty, strictly increasing enumeration of exactly those
zero-based rounds \(b\leq q\) for which the positive data slot
\(b.\mathrm{succ}\) is a PE with supporting vertex \(u\).  Then, for
every \(0\leq a<p\),
\[
 \gamma_{b_{a+1}+1}=\delta_{b_a+1}.
\]
In Lean the enumeration is an explicit map
\(\mathrm{Fin}(p+1)\to\mathrm{Fin}(r)\); strict increase and the
``if and only if'' description of its range are separate hypotheses.
\end{lemma}
\begin{proof}
Fix \(a<p\), and write \(b=b_a\) and \(b'=b_{a+1}\).  The exact-range
hypothesis says that both positive slots \(b.\mathrm{succ}\) and
\(b'.\mathrm{succ}\) have extension kind PE and chosen vertex \(u\).
Use the certificate field of \noderef{TashkinovSeries} at the latter
slot.  The kind equality eliminates the SE and RE constructors of
\noderef{TashkinovSeriesStep}, leaving its PE constructor.

Let \(g\) be the constructor's chosen value of \(\gamma_{b'+1}\).
The constructor records the history-sensitive rule
\noderef{ParallelExtensionColor}.  There is already an earlier PE at
\(u\), namely the slot \(b.\mathrm{succ}\), and strict increase gives
\(b<b'\).  Hence the ``no earlier PE at \(u\)'' alternative in that
rule is impossible.  The other alternative supplies a positive slot
\(h\), at most the preceding-coloring index \(b'\), such that
\(h\) is a PE at \(u\),
\[
 \delta_h=g,                                                     \tag{1}
\]
and every earlier PE slot at \(u\) through that index is at most
\(h\).

Translate the positive slot \(h\) to its zero-based round
\(d=h-1\).  Positivity makes this translation exact.  Since
\(d<b'\leq q\), the exact-range hypothesis places \(d\) in the displayed
enumeration.  Maximality in
\noderef{ParallelExtensionColor}, applied to
\(b.\mathrm{succ}\), gives \(b\leq d\).  There is no enumerated index
strictly between the consecutive entries \(b_a\) and \(b_{a+1}\): if
\(d=b_t\), strict increase would force \(t\leq a\) from \(d<b_{a+1}\),
whereas \(b_a\leq d\) forces \(a\leq t\).  Thus \(t=a\) and
\(d=b\), so \(h=b.\mathrm{succ}\).

Finally, the PE constructor identifies the optional gamma field at the
new slot with \(g\), while (1) identifies the optional delta field at the
preceding enumerated slot with the same \(g\).  Therefore
\(\gamma_{b_{a+1}+1}=\delta_{b_a+1}\), including equality of the exact
stored color values.  Since \(a<p\) was arbitrary, all successive
equalities follow.  The condition
\(u\in V(T_{q+1})\), expressed by
\noderef{TreeSequenceVertexSet}, is the manuscript's ambient condition
on the common supporting vertex.
\end{proof}
